\documentclass[12pt letter]{report}
\input{./template/preamble}
\input{./template/macros}
\input{./template/letterfonts}

\title{\Huge{The Link Layer and LANs}}
\author{\huge{Madiba Hudson-Quansah}}
\date{}
\usepackage{parskip}
\usepackage{tabularx}

\setcounter{tocdepth}{4}
\setcounter{secnumdepth}{4}

\begin{document}
\maketitle
\newpage
\pdfbookmark[section]{\contentsname}{too}
\tableofcontents
\pagebreak

\chapter{Introduction}

In link layer terminology, any device that runs a link-layer
protocol, i.e. layer 2 protocol is a \textbf{node}. Nodes include
hosts, routers, switches and WiFi access points.

The communication channels that connect adjacent nodes along the
communication path are called \textbf{links}. In order for a datagram
to be transferred from source node to destination node, it must be
moved over each of the individual links in the end-to-end path.

Over a given link a transmitting node encapsulates the datagram in a
link-layer frame and transmits the frame onto the link.

\section{Services provided by the Link Layer}

The possible services that can be offered by a link-layer protocol include:
\begin{description}
  \item[Framing]  - Encapsulating network-layer datagrams within a
    link-layer frame before transmission over the link. A frame
    consists of a data field, where the network-layer datagram is
    inserted,  and a number of header fields, with the structure
    specified by the link-layer protocol.
  \item[Link Access] - A medium access control (MAC) protocol
    specifies the rules by which a frame is transmitted onto the
    link. For point-to-point links that have a single sender and
    receiver the MAC protocol is trivial/non-existent. For shared
    links the MAC protocol coordinates the frame transmission by multiple nodes.
  \item[Reliable Delivery] - Guarantees to move each network-layer
    datagram across the link without error. This is used for links
    that are prone to high error rates, such as wireless links, with
    the goal of correcting the error locally, rather than forcing an
    end-to-end retransmission.
  \item[Error Dectection and Correction ] - Can detect and possibly
    correct errors that occur when a frame is
    transmitted over the link. This is used for links that are
    prone to high error rates, such as wireless links.
\end{description}

\section{Link Layer Implementation}

For the most part link-layer functionality is implemented on a chip
called the \textbf{Network Adapter} or \textbf{Network Interface
Card/Controller (NIC)}. This card implemented services including
framing, link access, error detection, etc.

On the sending side the NIC takes a datagram that has been created
and stored in host memory by the higher layers of the protocol stack,
encapsulates it in a link-layer frame, and then transmits the frame
into the communication link.

On the receiving side the NIC receives the entire frame and extracts
the network-layer datagram. If the link-layer performs error
detection, then it is the sending NIC that sets the error-detection
bits in the frame header, and the receiving NIC that checks these
bits to determine if an error has occurred.

\chapter{Error Detection and Correction Techniques}

Considering a node sending data $D$ over a link. To be protected
against but errors the data $D$ is augmented with error-detection and
correction bits $EDC$. Both $D$ and $EDC$ are sent to a receiving
node, where at the recipient a sequence of bits $D^{\prime}$ and
$EDC^{\prime}$ is received, which may differ from the original $D$
and $EDC$ due to errors that may have occurred during transmission.

The receiver's challenge is to determine whether or not $D^{\prime}$
is the same as $D$, given that it has only received $D^{\prime}$ and
$EDC^{\prime}$. Error-detection and correction techniques allow the
receiver to sometimes detect that bit errors have occurred, however
there may still be \textbf{undetected but errors}, where the receiver
may be unaware that $D^{\prime}$ contains bit errors. As a result the
receiver may deliver a corrupted datagram to the network-layer or be
unaware that the contents of field in the frame's header has been corrupted.

\section{Parity Checks}

The use of a \textbf{parity bit} is a simple error-detection
technique that can be used to detect single-bit errors.

Given data $D$ containing $d$ bits, in an \textbf{even parity}
scheme, the sender includes an additional but and chooses its value
such that the total number of 1s in the $d + 1$ bits is even. For
\textbf{odd parity} schemes, the parity but is chosen such that the
total number of 1s in the $d + 1$ bits is odd.

On receipt the receiver counts the number of 1s in the received $d +
1$ bits. If an odd number of of 1s are received in an even parity
scheme, or an even number of 1s are received in an odd parity scheme,
then the receiver can conclude that a single-bit error has occurred.

This is a simple to implement approach however measurements have
shown that rather than occurring independently, errors are often
clustered together in bursts. Under these conditions the probability
of undetected errors in a parity scheme can approach 50\%.

This parity scheme can however be generalized to use multiple parity
bits which forms the basis of other error-detection and correction
techniques, such as the Hamming code.

Again given there are $d$ bits in $D$ which are divided into $i$ and
$j$ rows and columns, a parity value is computed for each row and
column. The resulting $i + j + 1$ parity bits comprise the link-layer
frame's error-detection bits as shown in the figure below.

\begin{figure}[H]
  \centering
  \[
    \begin{array}{cccc|c}
      d_{1,1} & \cdots & d_{1,j} & & d_{1,j+1} \\
      d_{2,1} & \cdots & d_{2,j} & & d_{2,j+1} \\
      \vdots  & \ddots & \vdots  & & \vdots    \\
      d_{i,1} & \cdots & d_{i,j} & & d_{i,j+1} \\ \hline
      d_{i+1,1} & \cdots & d_{i+1,j} & & d_{i+1,j+1}
    \end{array}
  \]
  \caption{Two-dimensional parity arrangement with row and column parity bits}
  \label{fig:parity}
\end{figure}

Suppose now that a single bit error occurs in the original $d$ bits
of information, with this parity scheme the parity of both the row
and column containing the flipped but will be in error. This allows
the receiver to both detect and correct the error using the indices
of the row and column parity bits that are in error.

\begin{figure}[H]
  \centering
  \[
    \begin{array}{ccccc|c}
      1 & 0 & 1 & 0 & 1 & 1\\
      1 & 1 & 1 & 1 & 0 & 0\\
      0 & 1 & 1 & 1 & 0& 1 \\
      \hline
      0 & 0 & 1 & 0 & 1 & 0
    \end{array}
  \]
  \caption{No errors}
\end{figure}

\begin{figure}[H]
  \centering
  \[
    \begin{array}{ccccc|c}
      1 & 0 & 1 & 0 & 1 & 1\\
      1 & 0 & 1 & 1 & 0 & 0\\
      0 & 1 & 1 & 1 & 0& 1 \\
      \hline
      0 & 0 & 1 & 0 & 1 & 0
    \end{array}
  \]
  \caption{Correctable single-bit error in row 2, column 2}
\end{figure}

\dfn{Forward Error Correction (FEC)}{
  The ability of the receiver to both detect and correct errors.
}

\section{Checksumming}

In checksumming the $d$ bits of data are treated as a sequence if
$k$-bit integers. A simple checksum can be computed by summing these
$k$-but integers and then using the resulting sum as the
error-detection bits. This allows the receiver to check the checksum
by taking the 1s complement of the sum of the received data,
including the checksum, and checking whether the result is all 0s.
This approach can be used to detect burst errors, however it is not
as robust as cyclic redundancy checks.

\section{Cyclic Redundancy Check (CRC)}

An error-detection technique based on cyclic redundancy check (CRC)
codes / polynomial codes.

Considering the $d$ bits of data $D$, the sender and receiver must
first agree on a $r + 1$ bit pattern known as the \textbf{generator},
$G$. We will require that the most significant bit of $G$ be 1, with
the key idea of CRC codes being that for the given $D$ the sender
will choose $r$ addition bits $R$ and append then to $D$ such that
the resulting $d + r$ bit pattern when interpreted as a binary number
is exactly divisible by the generator $G$ using modulo-2 arithmetic.

The error checking process then involves the receiver dividing the
received $d + r$ bits by $G$, if the remainder is nonzero an error
has occurred, else the receiver can conclude that no errors have occurred, i.e.:
\[
  \left(D \cdot 2^r \oplus  R \right) \mod G = 0
\]

CRC calculations being done in modulo-2 arithmetic means that
addition and subtraction are performed using the exclusive OR (XOR),
thus the division process can be implemented using only XOR
operations, which are simple to implement in hardware.

The sender computes using the following procedure. The sender wants
to find $R$ such that there is an $n$ such that
\[
  D \cdot 2^r \oplus R = nG
\]
Making $R$ the subject we have
\begin{align*}
  D \cdot 2^r \oplus R &= nG\\
  D\cdot 2^r &= nG \oplus R \\
  R &= \text{remainder } \frac{D\cdot 2^r}{G} \\
\end{align*}

\ex{}{
  \qs{}{
    Given $D = 101110$, $d = 6$, $G = 1001$, and $r = 3$, compute the
    value of $R$ and the resulting $d + r$ bit pattern that is transmitted.
  }

  \sol{
    \begin{align*}
      R &= \text{remainder } \frac{D\cdot 2^r}{G} \\
      &= \text{remainder } \frac{101110 \cdot 2^{3}}{1001} \\
      &= \text{remainder } \frac{101110000}{1001} \\
      &= 011 \\
      \therefore D \cdot 2^r \oplus R &= 101110000 \oplus 011 \\
    \end{align*}
    And the transmitted $d + r$ bit pattern is $101110011$.
  }
}

\chapter{Multiple Access Links and Protocols}

There are two types of network links:
\begin{description}
  \item[Point-to-Point] - Single sender and receiver, thus no need
    for a medium access control
    protocol.
  \item[Broadcast Link] - Multiple senders and receivers, thus a
    medium access control protocol is needed to coordinate the access
    to the shared communication medium.
\end{description}

The term broadcast is used to refer to the fact that when any one
node transmits a frame, the channel broadcasts the frame and each of
the other nodes on the link receives a copy.

The coordination of access of multiple sending and receiving nodes
defines the \textbf{multiple access problem}. \textbf{Multiple Access
Protocols} are protocols by which nodes regulate their transmission
into the shared broadcast channel to resolve the multiple access problem.

Because all nodes are capable of transmitting frames, more than two
nodes can transmit frames at the same time, when this happens al of
the nodes receive multiple frames at the same time, i.e. the
transmitted frames collide at all of the receivers. Typically in the
event of a collision, none of the receiving nodes can  make any sense
of any of the frames that were transmitted, thus all the frames
involved in the collision are lost and the broadcast channel is
wasted during the collision interval.

In order to ensure that the broadcast channel performs useful work
when multiple nodes are active is it necessary to coordinate the
transmissions of the active nodes, thus the need for multiple access protocols.

Ideally, a multiple access protocol for a broadcast channel of rate
$R$ bits per second should have the following characteristics:

\begin{enumerate}
  \item  When only one node has data to send that node has a
    throughput of $R$ bps.
  \item When $M$ nodes have data to send each of the nodes has a
    throughput of $R/M$ bps on average.
  \item The protocol is decentralized, i.e. there is no master node
    that represents a single point of failure for the network.
  \item The protocol is simple, so that is inexpensive to implement.
\end{enumerate}

\section{Channel Partitioning Protocols}

Time-Division Multiplexing (TDM) and Frequency-Division Multiplexing
(FDM) are examples of channel
partitioning protocols, where the channel is divided into time slots
or frequency bands respectively, and each node is allocated a
specific time slot or frequency band for transmission.

TDM divides time into \textbf{time frames} and further divides each
frame into $N$ \textbf{time slots}. Each time slot is then assigned
to one of the $N$ nodes. Whenever a node has a packet to send, it
transmits the packet's bits during its assigned time slot in the
revolving TDM frame. TDM is appealing because it eliminates
collisions and is perfectly fair, which each node receiving a
dedicated transmission rate of $\frac{R}{N}$ bps during each frame time.
However, it has two major drawbacks:
\begin{itemize}
  \item A node is limited to an average of $\frac{R}{N}$ bps even
    when it is the only node with packets to send.
  \item A node must always wait for its turn in the transmission
    sequence, even if the other nodes have no packets to send.
\end{itemize}

FDM divides the $R$ bps channel into $N$ different frequencies, each
with a bandwidth of $\frac{R}{N}$ and assigns each frequency to one
of $N$ nodes. FDM thus creates $N$ smaller channels of $\frac{R}{N}$
bps out of the single larger $R$ bps channel. FDM shares the same
advantages and disadvantages as TDM, however it is more expensive to
implement than TDM, as it requires more complex hardware to generate
and filter the
different frequencies.

Another channel partitioning protocol is Code Division Multiple
Access (CDMA). CDMA assigns a different code to each node, and then
uses this code to encode the data bits it sends. If the codes are
chosen to be orthogonal, then multiple nodes can transmit
simultaneously without interference, as the receiver can use the code
to decode the received signal and extract the data bits for each node.

\section{Random Access Protocols}

Random access protocols allow nodes to transmit at the full channel
rate $R$ bps when they have data to send. When there is a collision
each node involved in the collision repeatedly retransmits its frame
until its frame gets through without collision, however this
retransmission doesn't happen immediately after the collision but
rather after a random delay independent of the other nodes, thus
reducing the probability of repeated collisions.

\subsection{Slotted ALOHA}

The following assumptions are made in slotted ALOHA:
\begin{itemize}
  \item All frames consist of exactly $L$ bits.
  \item Time is divided into slots of size $\frac{L}{R}$ seconds, i.e.
    a slot equals the time to transmit one frame.
  \item  Nodes start to transmit frames only at the beginning of slots.
  \item Nodes are synchronized so that each node knows when the slots begin.
  \item If two or more frames collide in a slot, then all the nodes
    detect the collision event before the slot ends.
\end{itemize}

Let $p$ be the probability that a node transmits in a slot, the
operation of slotted ALOHA can be described as follows:
\begin{itemize}
  \item When the node has a fresh frame to send, it waits until the
    beginning of the next slot and transmits the entire frame in the slot.
  \item If there isn't a collision, the node has successfully
    transmitted its frame and thus need not consider retransmitting the frame.
  \item If there is a collision, the nodes detects the collision
    before the end of the slot. The node retransmits its frame in
    each subsequent slot with probability $p$ until the frame is
    successfully transmitted without collision.
\end{itemize}

By retransmitting with probability $p$ rather than retransmitting
immediately after a collision, the probability of repeated collisions
is reduced, thus improving the efficiency of the protocol.

Slotted ALOHA has the following advantages:
\begin{itemize}
  \item Allows a node to transmit continuously at the full rate $R$,
    when it is the only active node
  \item  It is highly decentralized, as there is no master node that
    controls the
    transmission of the other nodes.
\end{itemize}

And the following disadvantages:
\begin{itemize}
  \item Requires slots to be synchronized, which can be difficult to
    achieve in practice.
  \item Collisions still occur, thus the channel is not used to its full
    capacity, even when there are only a few active nodes.
  \item When the number of active nodes is large, the efficiency of the
    protocol degrades significantly, as the probability of collisions
    increases, thus the channel is used very inefficiently.
\end{itemize}

In the case of multiple active nodes a certain fraction of the slots
will have collision and will therefore be wasted. Also another
fraction of the slots will be empty because all active nodes refrain
from transmitting as a result of the probabilistic transmission
policy. The only unwasted slots will be those in which exactly one
node transmits. This is called a \textbf{successful slot}.

The efficiency of a slotted multiple access protocol is defined as
the long-run fraction of successful slots in the case when there are
a large number of active nodes, each always having a large number of
frames to send.

The maximum efficiency of slotted ALOHA is $\frac{1}{e} \approx
0.37$, which occurs when $p = \frac{1}{M}$, where $M$ is the number
of active nodes. This means that even under ideal conditions, slotted
ALOHA can only
utilize about 37\% of the channel capacity, thus it is not a very
efficient protocol, for example given a 100-Mbps channel, the maximum
throughput of slotted ALOHA is only about 37 Mbps.

\subsection{ALOHA}

ALOHA, also called pure ALOHA, is an earlier variant of slotted ALOHA
where nodes are not required to synchronize their transmissions to
slot boundaries. In ALOHA a node transmits a frame whenever it has a
frame to send, and if there is a collision the node waits for a
random amount of time before retransmitting the frame.

The maximum efficiency of ALOHA is $\frac{1}{2e} \approx 0.18$, which
is significantly lower than the maximum efficiency of slotted ALOHA,
thus the requirement for slot synchronization
in slotted ALOHA is justified by the significant improvement in efficiency.

\subsection{Carrier Sense Multiple Access (CSMA)}

Carrier Sense Multiple Access (CSMA) is a random access protocol
in which a node with a frame to send first senses the channel to
determine whether it is idle or busy. If the channel is idle, the
node may transmit immediately (depending on the CSMA variant). If
the channel is busy, the node defers transmission according to the
protocol’s back off strategy.

Collisions can still occur in CSMA due to non-zero propagation delay.
Specifically, there is a delay between when one node begins transmitting
and when other nodes can detect that the channel is busy. As a result,
another node may sense the channel as idle and begin transmitting before
the first signal reaches it, leading to a collision.

\subsection{Carrier Sense Multiple Access with Collision Detection (CSMA/CD)}

CSMA with Collision Detection improves upon CSMA by allowing nodes to
detect collisions while they are transmitting. When a node performs
collision detection it stops transmission as soon as it detects a collision.

The operation of CSMA/CD can be described as follows:
\begin{enumerate}
  \item The node obtains a datagram from the network layer, prepares
    a link-layer frame and puts the frame in its buffer.
  \item If the node senses that the channel is idle, it starts to
    transmit the frame. If the node senses that the channel is busy,
    it defers transmission until the channel becomes idle.
  \item While transmitting the node monitors for the presence of an
    active signal on the channel coming from other active nodes.
  \item If the node transmits the entire frame without detecting
    activity from other nodes the node is finished with the frame.
    Else if the node detects activity from other nodes, it aborts transmission.
  \item After aborting the node waits a random amount of time and
    returns to step 2 to try again.
\end{enumerate}

The random back off time is determined by the \textbf{exponential /
binary exponential back off} algorithm, which works increases the
back off time exponentially with the number of collisions that have
occurred for the frame. Specifically, after $i$ collisions the node
chooses a random back off time uniformly from the interval $[0, 2^i - 1]$.

\subsubsection{CSMA/CD Efficiency}

The efficiency of CSMA/CD can be approximated by the following formula:
\[
  \text{Efficiency} = \frac{1}{1 + 5d_{\text{prop}} / d_{\text{trans}}}
\]

Where $d_{\text{prop}}$ is the maximum propagation delay between any
two nodes in the network, and $d_{\text{trans}}$ is the time to
transmit a frame. This formula shows that the efficiency of CSMA/CD
decreases as the ratio of propagation delay to transmission time increases.

\section{Taking Turns Protocols}

Taking turns protocols are multiple access protocols designed to
achieve an average throughput of $\frac{R}{M}$ bps when $M$ nodes
have data to send, while also allowing a node to transmit at the full
channel rate $R$ bps when it is the only active node.

\subsection{Polling}

The Polling Protocol requires one of the nodes to be designated the
master node, which polls each of the nodes in a round-robin fashion
to determine if they have data to send. If a node has data to send,
it is allowed to transmit its frame immediately after being polled by
the master node.

The polling protocol eliminates collision and empty slots allowing it
to achieve a higher efficiency. However, it has the following disadvantages:
\begin{itemize}
  \item There is a polling delay, i.e. the amount of time required to
    notify a node that it can transmit, which can be significant when
    the number of nodes is large.
  \item The master node represents a single point of failure for the
    network, as if the master node fails, the entire network will fail.
\end{itemize}

\subsection{Token-Passing}

The Token-Passing Protocol eliminates the need for a master node by
having nodes exchange a small, special-purpose frame known as a
\textbf{token} in a fixed order. When a node receives a token it
holds onto it only if it has some frames to transmit, else it
immediately passes the token to the next node.

Token passing is decentralized an highly efficient, however it has
the following disadvantages:
\begin{itemize}
  \item The failure of one node can crash the entire channel, as the
    token can be lost or the token circulation can be disrupted.
  \item If a node neglects to release the token after transmitting,
    the entire channel can be disrupted and some recovery procedure
    must be implemented to restore the token and resume normal operation.
\end{itemize}

\chapter{Switched Local Area Networks}

\section{Link-Layer Addressing and Address Resolution Protocol (ARP)}

Hosts and routers have link-layer addresses, which are used by the
link-layer to identify
the source and destination of a frame. Link-layer addresses are
attached to \textbf{network interfaces} and not to hosts or routers.

\subsection{Media Access Control (MAC) Addresses}

A host or router with multiple network interfaces will have multiple
link-layer addresses for each of its interfaces. However link-layer
switches do not have link-layer addresses associated with their
interfaces that connect hosts and routers, because the job of the
link-layer switch is to forward datagrams between hosts and routers,
which it does transparently, i.e. without the host or router having
to explicitly address the frame to the intervening switch.

A link-layer address is also called a \textbf{LAN address / Physical
address / Media Access Control (MAC) address}. The MAC address is 6
bytes long giving $2^{48}$ possible MAC addresses which is commonly
represented in hexadecimal notation as 12 hexadecimal digits.

No two network interfaces have the same address as the IEEE manages
the MAC address space, requiring a company to purchase a chunk of
address space consisting of $2^{24}$ addresses. An interface's MAC
address has a flat structure and doesn't change no matter where the
interface is connected in the network, thus it is not hierarchical
like IP addresses. A MAC address is analogous to a person's social
security number, while an IP address is analogous to a person's home address.

When an interface wants to send a frame to some destination
interface, the sending interface inserts the destination interface's
MAC address into the frame and sends the frame into the LAN. Due to
the broadcast nature of the LAN, all interfaces on the LAN receive a
copy of the frame, however only the destination interface will
process the frame, while the other interfaces will simply discard the
frame. However sometimes a sending node wants all nodes on the LAN to
receive and process a frame, in this case the sending node specifies
a special MAC \textbf{broadcast address}, i.e FF:FF:FF:FF:FF:FF, as
the destination address in the frame, which causes all nodes on the
LAN to process the frame.

\subsection{Address Resolution Protocol (ARP)}

Because there are both network-layer and link-layer addresses, there
is a need to translate between them, this is done using the Address
Resolution Protocol (ARP).

ARP resolves the MAC address of a destination host given its IP
address. Suppose a host with IP address 222.222.222.220 wants to send
an IP datagram to a host with IP address 222.222.222.222 in the same
subnet. To send this datagram the source must give its NIC not only
the IP datagram but also the MAC address for the destination host. An
ARP module in the sending host takes any IP address on the same LAN
as input and returns the corresponding MAC address as input.

Each host and route has an \textbf{ARP table}, which contains
mappings of IP address to MAC addresses. The ARP table also contains
a \textbf{time-to-live (TTL)} which indicates when each mapping will
be deleted from the table.

Going back to the earlier example supposed the sending host wants to
send a datagram that is IP addressed to another host or router on the
same subnet. IN the case where the sending host's ARP table contains
a valid mapping for the destination IP address, the sending host can
simply use the corresponding MAC address to send the datagram.
However, if the sending host's ARP table doesn't contain a valid
mapping for the destination host the following process occurs:
\begin{enumerate}
  \item First the sender constructs an ARP packet, which contains
    several fields including the sending and receiving IP and MAC
    addresses. Both ARP query and response packets have the same
    format. The purpose of the ARP query packet is to query all the
    other hosts and routers on the subnet to determine the MAC
    address corresponding to the destination IP address, while the
    purpose of the ARP response packet is to respond to an ARP query
    with the requested MAC address.
  \item The sender passes the ARP query packet to its NIC along with
    an indication that the packet should be broadcast to all nodes on
    the subnet using the broadcast MAC address.
  \item The ARP query is received by all the other nodes on the
    subnet and each node checks to see if its IP address matches the
    destination IP address in the ARP packet. The node that does
    match sends back to the querying hot a response ARP packet with
    the desired mapping.
  \item The querying host then updates its ARP table and sends its IP
    datagram to the destination host using the newly acquired MAC address.
\end{enumerate}

It is important to note that the ARP query message is within a
broadcast frame, whereas the response ARP message is sent within a
standard frame. This is because at the time of the ARP query the
sender doesn't know the MAC address of the destination host, thus it
must use the broadcast MAC address to ensure that all nodes on the
subnet receive the ARP query, whereas at the time of the ARP response
the sender knows the MAC address of the destination host, thus it can
use this MAC address to send the ARP response directly to the destination host.

ARP is also "plug-and-play", i.e. an ARP table gets built
automatically as hosts and routers send ARP queries and responses,
thus there is no need for a network administrator to manually
configure the ARP tables.

ARP is in a grey area between a link-layer and network-layer
protocol, as it is used by the link-layer to resolve IP addresses to
MAC addresses, however it is implemented as a network-layer protocol,
as it uses IP datagrams to carry ARP query and response messages,
thus it is implemented in software and runs on the host's CPU rather
than in the NIC hardware.

\subsection{Sending a Datagram off the Subnet}

\dfn{Hop}{
  A hop is a single step in the path from a source host to a
  destination host, i.e. the transmission of a datagram from one node
  to an adjacent node. Denoted by a router in the path, as a datagram
  is forwarded from one router to the next until it reaches the
  destination host.
}

When a host wants to send a datagram to a host that isn't in its
subnet, the sending host must send the datagram to a router on the
same subnet as the sending host, which will then forward the datagram
to the destination host.

The destination IP address of the datagram is still the IP address of
the final destination host and does not change while traversing the
network, however the destination MAC address changes at each hop as
the datagram is forwarded from the sending host to the first-hop
router, and then from the first-hop router to the second-hop router,
and so on until the datagram reaches the destination host.

A router has an IP address for each of its interfaces, and for each
interface there is a corresponding ARP module and NIC. Given a host A
with IP address 111.111.111.111 and MAC address AA:AA:AA:AA:AA:AA,
and a router R with two interfaces, one with IP address
111.111.111.110 and MAC address RR:RR:RR:RR:RR:RR, and the other with
IP address 222.222.222.220 and MAC address RR:RR:RR:RR:RR:RS, when
host A wants to send a datagram to host B with IP address
222.222.222.221 with MAC address BB:BB:BB:BB:BB:BB, the following
process occurs:
\begin{enumerate}
  \item Host A creates a datagram with network-layer destination
    address 222.222.222.221. It then passes the datagram to its NIC
    but instead of broadcasting the datagram, host A sends the
    datagram to router R's interface with IP address 111.111.111.110,
    i.e. the IP address of the first-hop router on the path to the final
    destination, with corresponding MAC address RR:RR:RR:RR:RR:RR.
    The host acquires the MAC address of the first-hop router using
    ARP if it doesn't have it using ARP within the subnet as
    described in the previous section.
  \item Router R receives this frame and processes it as it is
    addressed to one of its interfaces. Router R now has to determine
    the correct interface on which the datagram should be forwarded.
    It does this through forwarding as described in the network layer
    chapter. Once the correct interface is determined, router R then
    needs to determine the MAC address of the destination host B,
    which it does using ARP within the subnet as described in the
    previous section. Finally, router R sends the datagram to host B
    using host B MAC address BB:BB:BB:BB:BB:BB as the destination MAC
    address in the
    frame.
  \item Host B receives the frame and processes it as it is addressed
    to one of its interfaces. Host B then extracts the datagram and
    delivers it to the network layer.
\end{enumerate}

\section{Ethernet}

\dfn{Hub}{
  A physical-layer device that acts on individual bits rather than frames.
}

Ethernet is the most widely used link-layer protocol for LANs. When a
bit arrives from one interface the \textbf{hub}  recreates the bit,
boosts its energy strength and transmits the bit onto all the other
interfaces. Ethernet with a hub-base star topology is also a
broadcast LAN, i.e. all nodes receives a copy of all frames
transmitted across the network. This means collisions can occur when
two or more nodes transmit at the same time, needing a retransmission
of the frames involved in the collision.

\subsection{Ethernet Frame Structure}

An Ethernet frame consists of the following fields:
\begin{description}
  \item[Data]  - 46 - 1500 bytes. Carries the IP datagram. The
    maximum transmission unit (MTU) of Ethernet is 1500 bytes, which
    means that if the datagram exceeds 1500 bytes the host has to
    fragment the datagram.
  \item[Destination MAC Address] - 6 bytes. Contains the MAC address
    of the destination node.
  \item[Source MAC Address] - 6 bytes. Contains the MAC address of
    the source node.
  \item[Type] - 2 bytes. Permits the Ethernet to multiplex
    network-layer protocols by indicating the type of the
    network-layer datagram being carried in the data field. This is
    needed because hosts can use other network-layer protocols
    besides IP, such as ARP, thus the type field allows the Ethernet
    to indicate which network-layer protocol is being used in the
    data field. This allows demultiplexing of the data field to the
    correct network-layer protocol at the receiving node.
  \item[Cyclic Redundancy Check (CRC)] - 4 bytes. Used for error detection.
  \item[Preamble] - 8 bytes. Used to synchronize the receiver's clock
    to the sender's clock. Each of the first 7 bytes of the preamble
    has a vale of 10101010, while the last byte has a value of
    10101011, which allows the receiver to detect the start of the
    frame and synchronize its clock to the sender's clock.
    Synchronization is needed because the sender and receiver may
    have different clock rates, thus the preamble allows the receiver
    to adjust its clock rate to match that of the sender.
\end{description}

Ethernet technologies provide:
\begin{itemize}
  \item Connectionless service
  \item Unreliable service
  \item Best effort delivery
\end{itemize}

\subsection{Ethernet Technologies}

\dfn{Repeater}{
  A physical-layer device that receives a signal on the input side
  and regenerates the signal on the output side.
}

Ethernet comes in many different varieties, which differ in the
physical medium they use, the maximum length of the medium, and the
data rate they support.

Initially Ethernet was designed to run over a segment of coaxial
cable, however it has evolved to run over twisted pair and optical
fiber cables. The original Ethernet standard, known as 10BASE5,
supported a data rate of 10 Mbps and a maximum segment length of 500
meters. The 10BASE2 standard also supported a data rate of 10 Mbps
but had a maximum segment length of 200 meters. The 10BASE-T
standard, which runs over twisted pair cables, also supports a data
rate of 10 Mbps but has a maximum segment length of 100 meters.

The names of Ethernet standards typically follow the format $X$BASE$Y$, where:
\begin{description}
  \item[$X$] - The data rate in Mbps, e.g. 10 for 10 Mbps, 100 for 100
    Mbps, 1000 for 1 Gbps, etc.
  \item[BASE] - Indicates that the standard uses baseband signalling, i.e.
    the entire bandwidth of the medium is used to transmit a single signal.
  \item[$Y$] - Indicates the type of physical medium used, e.g.
    T for twisted pair, F for fiber optic, etc.
\end{description}

\section{Link-Layer Switches}

\dfn{Switch}{
  A link-layer device that acts on frames rather than bits, i.e. it
  receives an entire frame, processes the frame and then forwards the
  frame to the correct outgoing interface.
}

\dfn{Transparent}{
  A device is transparent to the hosts and routers connected to it if
  the hosts and routers are not aware of the presence of the device in
  the network, i.e. the device operates without the need for any
  configuration or management by the hosts and routers.
}

Switches are \textbf{transparent} to the hosts and routers connected to them.

\subsection{Forwarding and Filtering}

\dfn{Filtering}{
  The switch function that determines whether a frame should be
  forwarded to some interface or should be dropped.
}

\dfn{Forwarding}{
  The switch function that determines the interfaces to which a frame
  should be direct, and moves the from to those interfaces.
}

Switch filtering and forwarding are done with a \textbf{switch
table}. The switch table contains entries for some of the hots and
routers on a LAN. AN entry in the switch table contains:
\begin{enumerate}
  \item A MAC address
  \item The switch interface that leads toward that MAC address
  \item The time at which the entry was placed in the table.
\end{enumerate}

Supposing a frame with destination address D arrives at the switch on
interface $x$. The switch indexes its table with the MAC address D.
There are three possible cases:
\begin{description}
  \item[There is no entry for D] - The switch forwards copies of the
    frame to the output buffers preceding all interfaces expect for
    interface $x$, i.e. broadcasting.
  \item[There is an entry in the table for D associating it with $x$]
    - In this case the frame is coming from a LAN segment that
    contains interface D. Thus there is no need to forward the frame
    to any other interfaces, the switch performs the filtering action.
  \item[There is an entry in the table for D associating it with some
    other interface $y$] - The frame needs to be forwarded to the LAN
    segment attached to interface $y$. The switch performs its
    forwarding function by putting the frame in an output buffer that
    precedes interface $y$.
\end{description}

In this sense a switch is smarter than a hub, as a switch can filter
frames and only forward frames to the correct outgoing interface,
while a hub simply forwards all bits to all interfaces.

\subsection{Self-Learning}

A switch is a \textbf{self-learning} device, i.e. its switch table is
build automatically, dynamically, and autonomously, without any
intervention from a network administrator or configuration protocol.
This is accomplished as follows:
\begin{enumerate}
  \item The switch table starts empty.
  \item  For each incoming frame received on an interface, the switch
    stores in its table
    \begin{enumerate}
      \item The MAC address in the frame's source address field
      \item The interface from which the frame arrived
      \item The current time
    \end{enumerate}
    In this manner the switch learns the location of the source host
    or router of the frame, and thus the switch can use this
    information to forward future frames destined to that host or router.
  \item The switch deletes an address in the table if no frames are
    received with that address as the source address after some
    period of time, \textbf{ageing time}. In this manner if a PC is
    replaced by another PC the MAC address of the original PC will
    eventually be deleted from the switch table.
\end{enumerate}

Switches are plug-and-play devices because they require no manual
configuration or management, as they are self-learning devices that
build their switch tables automatically and dynamically.

\subsection{Properties of Link-Layer Switching}

Switches therefore have the following advantages over broadcast links
such as buses or hub-based star topologies:
\begin{description}
  \item[Elimination of collisions]  - There is no wasted bandwidth
    due to collision, as the switches buffer frames and never
    transmit more than one frame on a segment at once. As with a
    router the maximum aggregate throughput of a switch is the sum of
    all the switch interface rates.
  \item[Heterogeneous links] - Because a switch isolates one link
    from another the different  links in the LAN can operate at
    different speeds and can run over different media.
  \item[Management] - Switches eases network management as they can
    detect and disconnect misbehaving nodes, gather statistics on
    bandwidth usage, collision rates and traffic types, and make it
    easier to add and move nodes in the network.
\end{description}

\subsection{Switches vs. Routers}

Although a switch is also a store-and-forward packet switch, it
differs from a router in that it forwards using MAC addresses rather
than IP addresses. Where a router is a layer-3 packet switch a switch
is a layer-2 packet switch.

Even though they are different devices a network admin must often
choose between then when installing an interconnection device. The
choice between a switch and a router depends on the specific needs of
the network, such as the size of the network, the amount of traffic,
and the need for security and control. In general, switches are more
suitable for smaller networks with less traffic, while routers are
more suitable for larger networks with more traffic and a need for
greater control and security.

A switch has the following advantages over a router:
\begin{itemize}
  \item There are plug and play devices that require no
    configuration, while routers require
    configuration and management.
  \item  They can have relatively high filtering and forwarding
    rates, as they only have to process frames only up to layer 2
    whereas routers have to process packets up to layer 3, thus
    switches can be faster than routers.
\end{itemize}

However a switch has the following disadvantages compared to a router:
\begin{itemize}
  \item  The active topology of a switched network is restricted to a
    spanning tree to prevent the cycling of broadcast frames.
  \item Large switched networks require large ARP tables in hosts and
    routers and would generate substantial ARP traffic and processing.
  \item Switches are also susceptible to broadcast storms, i.e. if
    one host misbehaves and transmits an endless stream of Ethernet
    broadcast frames, the switch will forward all these frames
    causing the entire network to collapse.
\end{itemize}

A router has the following advantages over a switch:
\begin{itemize}
  \item Packets do not normally cycle through routers even with
    redundant paths, as network addressing is often hierarchical, and
    thus packets are not restricted to a spanning tree and can use
    the best path between source and destination.
  \item Routers provide firewall protection against layer-2 broadcast storms.
\end{itemize}

However a router has the following disadvantages compared to a switch:
\begin{itemize}
  \item Routers are not plug-and-play devices, as they require configuration and
    management by a network administrator, while switches are
    plug-and-play devices that require no configuration or management.
\end{itemize}

\section{Virtual Local Area Networks (VLANs)}

Modern LANs are often configured hierarchically, which each workgroup
having its own switched LAN connected to the switched LANs of other
groups via a switch hierarchy. There are drawbacks to this configuration:
\begin{description}
  \item[Lack of traffic isolation]  - Because the LAN is shared by
    all the workgroups, traffic from one
    workgroup can affect the performance of other workgroups, e.g. if
    one workgroup generates a large amount of traffic it can cause
    congestion and delay for the other workgroups.
  \item[Inefficient use of switches] - Because the LAN is shared by
    all the workgroups, the switches in the LAN are shared by all the
    workgroups, thus if one workgroup generates a large amount of
    traffic it can cause congestion and delay for the other
    workgroups, and thus the switches in the LAN are not used efficiently.
  \item[Managing users] - If a user moves from one workgroup to
    another, the network administrator must reconfigure the switch
    ports to which the user is connected, which can be time-consuming
    and error-prone.
\end{description}

These drawbacks can be addressed by using Virtual Local Area Networks (VLANs).

\dfn{Virtual Local Area Network (VLAN)}{
  A VLAN is a logical LAN that is created by partitioning a physical
  LAN into multiple logical LANs. Each VLAN is a separate broadcast
  domain, i.e. frames sent by a host in one VLAN are not received by
  hosts in other VLANs, thus providing traffic isolation between
  different workgroups.
}

Switches that support VLANs allow multiple VLANs to be defined over a
single physical LAN infrastructure. Hosts within a VLAN communicate
with each other as if they were the only hosts connected to the switch.

\chapter{Link Virtualisation}

\section{Multiprotocol Label Switching (MPLS)}

Multiprotocol Label Switching (MPLS) is a link-layer protocol that
allows for the creation of virtual links between nodes in a network.

\chapter{Data Centre Networking}

\section{Data Centre Architectures}

\subsection{Load Balancing}

\subsection{Hierarchical Architecture}

\section{Trends in Data Centre Networking}

\subsection{Cost Reduction}

\subsection{Centralized SDN Control and Management}

\subsection{Virtualisation}

\subsection{Physical Constraints}

\subsection{Hardware Modularity and Customization}

\end{document}
